\documentclass[11pt]{article}

\usepackage[utf8]{inputenc}
\usepackage[T1]{fontenc}
\usepackage[margin=1in]{geometry}
\usepackage{amsmath,amssymb,amsfonts}
\usepackage{graphicx}
\usepackage{booktabs}
\usepackage{hyperref}
\usepackage{xcolor}
\usepackage{multirow}
\usepackage[numbers,sort&compress]{natbib}

\title{\textbf{Constraint Residual Detection: A Geometry-Based Approach to LLM Hallucination Detection}}

\author{
  Xinhang Deng\textsuperscript{1}\thanks{Correspondence: sampson1735937149@gmail.com. ORCID: 0009-0009-3974-8554.} \\
  \textsuperscript{1}Independent Research, Shenzhen, China \\
}

\date{May 2026}

\begin{document}
\maketitle

\begin{abstract}
Large language models (LLMs) produce hallucinations with high confidence, rendering uncertainty-based detection ineffective on a substantial fraction of factual errors. On a 30-question false-premise subset of TruthfulQA, Qwen2.5-7B-Instruct produces confident wrong answers on 36.7\% of questions---precisely the regime where predictive entropy and self-consistency methods fail (AUC $\leq$ 0.465). We propose \textbf{Constraint Residual Detection}, a method that measures violations of the model's internal constraint field $\Pi = \sum_i \nabla\sigma_i$ across the input-to-output transition. Rather than observing output confidence, we extract the constraint gradient field from hidden states and compute the norm change $\Delta\|\Pi\|$ as a hallucination signal. On TruthfulQA with Qwen2.5-7B-Instruct, our method achieves AUC = 0.816 (Cohen's $d$ = 1.06, $p$ = 0.018), compared to 0.465 for predictive entropy and 0.250 for SelfCheckGPT. Cross-scale analysis reveals that the signal weakens substantially at smaller model sizes: AUC = 0.532 at 1.5B parameters (hidden dim 1536) versus 0.816 at 7B (hidden dim 3584). The method adds $\sim$28.5ms of computational overhead beyond standard generation (total latency $\sim$1.1s).
\end{abstract}

\section{Introduction}

Hallucination---the generation of plausible but factually incorrect content---remains a critical barrier to deploying LLMs in production. The dominant detection paradigm measures output uncertainty: predictive entropy~\cite{kadavath2022language}, token probability, or multi-sample consistency~\cite{manakul2023selfcheckgpt,kuhn2023semantic}. These methods share a core assumption: incorrect outputs exhibit higher uncertainty than correct ones.

This assumption fails systematically on false-premise questions. When asked ``Can sharks get cancer?'', a model confidently responds ``No, sharks are immune to cancer''---and uncertainty-based detectors report low risk. The model \textit{believes} its hallucination, so surface confidence decouples from correctness. On our 30-question false-premise evaluation set (TruthfulQA~\cite{lin2022truthfulqa}), Qwen2.5-7B-Instruct hallucinates on 36.7\% of questions (11/30), and competing detectors perform near or below random (AUC $\leq$ 0.465).

We propose a different approach. Every LLM acquires implicit constraints during training---facts, logical relationships, syntactic structures. Content that violates these internalized constraints produces a detectable geometric signal in hidden state dynamics, \textit{regardless of surface confidence}. This connects to recent work showing that truth-relevant information is linearly accessible in LLM representations~\cite{burns2022discovering,azaria2023internal}, but differs in measuring \textit{constraint violation dynamics} across the input-to-output transition rather than probing static hidden states.

Mathematically, we define the \textbf{constraint field} $\Pi(p) = \sum_i \nabla\sigma_i(p)$ where each $\sigma_i$ measures a distinct dimension of internal constraint, and $\nabla\sigma_i$ captures token-level violation. The core empirical finding is:

\begin{equation}
\Delta\|\Pi\| = \|\Pi_{\text{output}}\| - \|\Pi_{\text{input}}\| \begin{cases}
> 0 & \text{(hallucination: constraint violation increases during generation)} \\
\leq 0 & \text{(factual: constraint violation is stable or resolves)}
\end{cases}
\label{eq:hypothesis}
\end{equation}

When the model hallucinates, it generates content that violates its internal constraints more than the input does, producing a positive $\Delta\|\Pi\|$. When it answers factually, constraint violation stays stable or decreases, yielding near-zero or negative $\Delta\|\Pi\|$. The detection signal is designed to be independent of surface confidence because it observes internal dynamics rather than output distributions.

\section{Method}

\subsection{Constraint Functions}

We operationalize five constraint functions $\sigma_i: \mathcal{H} \rightarrow [0,1]$, each capturing a distinct dimension of internal model behavior. All $\sigma_i$ are normalized to $[0,1]$ with NaN/Inf guards.

\paragraph{Factual Consistency ($\sigma_{\text{fact}}$).}
We operationalize factual consistency via the maximum attention weight at token position $t$ as a decision-certainty proxy: when the model is confident about a token---factual or not---attention concentrates on specific preceding positions:
\begin{equation}
\sigma_{\text{fact}}(h_t) = \max_{j \leq t} A_t[j]
\end{equation}
where $A_t \in \mathbb{R}^n$ is the attention weight vector at position $t$ (averaged across heads). In our experiments, attention weights are always extracted (\texttt{output\_attentions=True}), so this is the formulation actually evaluated. As an alternative, a truth direction $\mathbf{v}_{\text{truth}} \in \mathbb{R}^d$ can be calibrated from 9 contrastive prompt pairs (truthful vs.\ false completions), following evidence that truthfulness has a linear direction in activation space~\cite{burns2022discovering}. Under that formulation, $\sigma_{\text{fact}}(h_t) = (\cos(h_t, \mathbf{v}_{\text{truth}}) + 1)/2$, where the direction is the normalized difference between mean truthful and mean false hidden states from the final transformer layer. A third fallback, used only when neither attention weights nor calibration is available, is an L2-norm stability proxy: $\sigma_{\text{fact}}(h_t) = 1.0 - |\|h_t\| - \bar{n}|/\bar{n}$, where $\bar{n}$ is the empirical mean hidden state norm for the target model.

\paragraph{Syntactic Well-formedness ($\sigma_{\text{syntax}}$).}
Using attention weights $A_t \in \mathbb{R}^n$ at position $t$ (averaged across heads):
\begin{equation}
\sigma_{\text{syntax}}(A_t) = 1 - \frac{H(A_t)}{H_{\max}}, \quad H(A_t) = -\sum_{j=1}^{n} a_j \log a_j, \quad H_{\max} = \log n
\end{equation}
Low entropy (focused attention) indicates stronger syntactic constraint satisfaction.

\paragraph{Stylistic Coherence ($\sigma_{\text{style}}$).}
Given layer-wise hidden states $\{h_t^{(l)}\}_{l=1}^{L}$ at token $t$, we measure cross-layer norm stability via the coefficient of variation:
\begin{equation}
\sigma_{\text{style}}(\{h_t^{(l)}\}) = \exp\left(-3 \cdot \frac{\text{std}(\{\|h_t^{(l)}\|\})}{\text{mean}(\{\|h_t^{(l)}\|\})}\right)
\end{equation}
Erratic norm variation across layers indicates register shifts or tonal inconsistency. The decay factor 3 was set on a small validation set.

\paragraph{Safety Alignment ($\sigma_{\text{safety}}$).}
Analogous to $\sigma_{\text{fact}}$, a refusal direction $\mathbf{v}_{\text{refusal}} \in \mathbb{R}^d$ is calibrated from 4 contrastive prompt pairs (refusal vs.\ compliant responses):
\begin{equation}
\sigma_{\text{safety}}(h_t) = \frac{\cos(h_t, \mathbf{v}_{\text{refusal}}) + 1}{2}
\end{equation}
Without calibration, $\sigma_{\text{safety}}$ defaults to 0.3 (neutral). This constraint captures the model's safety-aligned operating regime, reflecting behaviors instilled during RLHF training~\cite{ouyang2022training}.

\paragraph{Semantic Coherence ($\sigma_{\text{coherence}}$).}
Adjacent-token cosine similarity:
\begin{equation}
\sigma_{\text{coherence}}(h_{t-1}, h_t) = \frac{\cos(h_{t-1}, h_t) + 1}{2}
\end{equation}
Sharp drops indicate semantic discontinuities that may signal confabulation.

\subsection{Constraint Gradient Field}

For each token position $t \geq 1$, per-constraint gradients are computed via finite differences:
\begin{equation}
\nabla\sigma_i(t) = \sigma_i(t) - \sigma_i(t-1)
\end{equation}
yielding a 5-dimensional gradient vector $\nabla\sigma(t) \in \mathbb{R}^5$.

The constraint residual at token $t$ is the signed sum:
\begin{equation}
\Pi(t) = \sum_{i=1}^{5} \nabla\sigma_i(t)
\end{equation}
Its magnitude $\|\Pi(t)\|$ captures the net constraint violation. We also define the cancellation ratio $c(t) = \|\sum \nabla\sigma_i\| / \sum \|\nabla\sigma_i\|$, which measures gradient component alignment: low $c(t)$ indicates opposing constraints, a signature of internal inconsistency. As a standalone detector, $c(t)$ performs near random (AUC = 0.528 at 1.5B), so we use it only as an auxiliary diagnostic rather than a primary detection signal.

\subsection{Detection Signal}

The detection signal is the mean residual magnitude shift from input to output:
\begin{equation}
\Delta\|\Pi\| = \frac{1}{|T_{\text{out}}|} \sum_{t \in T_{\text{out}}} \|\Pi(t)\| - \frac{1}{|T_{\text{in}}|} \sum_{t \in T_{\text{in}}} \|\Pi(t)\|
\label{eq:delta}
\end{equation}
where $T_{\text{in}}$ are content tokens in the input (excluding system prompts and special tokens), and $T_{\text{out}}$ are generated output tokens, with near-zero residuals ($\|\Pi(t)\| \leq 10^{-6}$) excluded to reduce noise. For hallucinations: $\Delta\|\Pi\| > 0$ (constraint violation increases). For factual answers: $\Delta\|\Pi\| \leq 0$ (constraint violation is stable or resolves).

In practice, hidden states for $T_{\text{in}}$ are obtained by encoding the input prompt in a separate forward pass (with \texttt{output\_hidden\_states=True}), and states for $T_{\text{out}}$ are obtained by re-encoding the generated output text independently. This means input constraint states are computed without autoregressive generation context, and output constraint states are computed without the prompt context. While this departs from the ideal of measuring constraint dynamics during generation, it avoids the engineering complexity of instrumenting the generation loop and adds negligible cost (see \S\ref{sec:latency}).

\subsection{Calibration}

To map the raw $\Delta\|\Pi\|$ signal to an interpretable probability, we fit a logistic regression on the evaluation data:
\begin{equation}
P(\text{hallucination} \mid \Delta\|\Pi\|) = \frac{1}{1 + e^{-(\alpha + \beta \cdot \Delta\|\Pi\|)}}
\label{eq:calibration}
\end{equation}

Fitting logistic regression (no regularization) on all $n=30$ evaluation samples yields $\alpha = -1.30$, $\beta = 16.52$. Under this calibration, $\Delta\|\Pi\| = -0.08$ maps to $P \approx 0.07$ and $\Delta\|\Pi\| = +0.08$ maps to $P \approx 0.51$. Since this is an in-sample fit on the evaluation data, the calibrated-probability AUC equals the raw-score AUC (0.816) by monotonicity. Two-fold cross-validation, a more conservative estimate, gives a mean test AUC of approximately 0.76. Risk tiers are defined as:
\begin{itemize}
  \item \textbf{Low}: $P < 0.32$ --- high confidence in correctness ($\Delta\|\Pi\| < 0.033$)
  \item \textbf{Medium}: $0.32 \leq P < 0.52$ --- uncertain, human review recommended ($0.033 \leq \Delta\|\Pi\| < 0.084$)
  \item \textbf{High}: $0.52 \leq P < 0.72$ --- likely hallucination ($0.084 \leq \Delta\|\Pi\| < 0.136$)
  \item \textbf{Critical}: $P \geq 0.72$ --- almost certain hallucination ($\Delta\|\Pi\| \geq 0.136$)
\end{itemize}

We emphasize that with $n = 30$, these are preliminary parameters. The 2-fold cross-validation indicates substantial estimation variance (fold-level AUC range: 0.72--0.80). Larger-sample calibration across diverse domains is needed for production-grade reliability.

\section{Experiments}

\subsection{Setup}

We evaluate on TruthfulQA~\cite{lin2022truthfulqa}, a benchmark of 817 questions probing common human misconceptions across 38 categories. We focus on the false-premise subset---questions containing false assumptions---which is the hardest case for uncertainty-based methods: the model's internalized misconceptions produce confident wrong answers.

\textbf{Model}: Qwen2.5-7B-Instruct (7B parameters, hidden dimension $d = 3584$, 28 layers, 28 query heads with grouped query attention, 4 KV heads).

\textbf{Evaluation}: $n = 30$ false-premise questions (11 hallucinatory, 19 correct; hallucination rate 36.7\%). All experiments run on a single NVIDIA GPU with 24GB VRAM. Hidden states are extracted via \texttt{output\_hidden\_states=True} on separate encoding passes (input and output), rather than during the autoregressive generation step. Hallucination labels are determined by self-judgment: the same model evaluates whether its own response contains factual errors given the reference answer, with a rule-based fallback on failure. This introduces a potential circularity concern---the model judges its own outputs---which we acknowledge as a limitation of the current labeling methodology.

\textbf{Baselines}:
\begin{itemize}
  \item \textbf{Predictive Entropy}: Token-level entropy of the output distribution~\cite{kadavath2022language}.
  \item \textbf{Max Probability}: Maximum token probability as confidence proxy.
  \item \textbf{SelfCheckGPT} (3$\times$ sampling): Multi-sample consistency via pairwise BERTScore~\cite{manakul2023selfcheckgpt}.
  \item \textbf{Response Length}: Heuristic based on response token count distributions.
\end{itemize}

Two additional methods were tested but omitted from the main comparison as they provided no discriminative signal on the false-premise subset: attention entropy (AUC = 0.500, at chance) and layer variance (AUC = 0.263, worse than chance). All methods share the same underlying model inference for fair latency comparison.

\subsection{Main Results}

\begin{table}[h]
\centering
\caption{TruthfulQA Hallucination Detection Benchmark (Qwen2.5-7B-Instruct, false-premise subset, $n = 30$)}
\label{tab:benchmark}
\begin{tabular}{lcccc}
\toprule
\textbf{Method} & \textbf{AUC} & \textbf{Precision} & \textbf{Recall} & \textbf{Latency (ms)} \\
\midrule
Constraint Residual (Ours) & \textbf{0.816} & 0.500 & 1.000 & 1,076 \\
Predictive Entropy & 0.465 & 0.167 & 0.333 & 1,005 \\
Max Probability & 0.447 & 0.167 & 0.333 & 1,006 \\
SelfCheckGPT (3$\times$) & 0.250 & 0.000 & 0.000 & 3,011 \\
Response Length & 0.443 & 0.250 & 0.500 & 1,004 \\
\midrule
\textit{Cohen's $d$ (hallucination vs.\ correct separation)} & \textit{1.06} & --- & --- & --- \\
\textit{$p$-value (independent-samples $t$-test, two-tailed)} & \textit{0.018} & --- & --- & --- \\
\bottomrule
\end{tabular}
\end{table}

Table~\ref{tab:benchmark} presents the results. Our method achieves AUC = 0.816, a +0.351 advantage over the best baseline (predictive entropy, AUC = 0.465). Our method achieves perfect recall (1.000) at 0.500 precision---detecting every hallucination while flagging half of the correct answers for review. This operating point was determined by optimizing the decision threshold on the evaluation set; out-of-sample recall may be lower. Nonetheless, the score distributions are well-separated (Cohen's $d$ = 1.06, $p$ = 0.018), confirming that the underlying signal is robust.

The baselines collectively fail to discriminate. Predictive entropy and max probability hover near random (AUC $\approx$ 0.45--0.47). SelfCheckGPT (AUC = 0.250) is \textit{below} random: it detects zero hallucinations (precision = recall = 0.000), assigning ``safe'' to all samples.

\subsection{Why Baselines Fail}

The SelfCheckGPT result (AUC = 0.250, precision = recall = 0) reveals a structural failure mode. On false-premise questions, the model has internalized the misconceptions being probed. Sampling three times yields three \textit{consistently wrong} answers, which the consistency check interprets as high agreement $\rightarrow$ ``safe.'' The method never flags a hallucination because consistency and correctness are inversely correlated on this task.

Predictive entropy and max probability fail because the model is well-calibrated to its own (incorrect) beliefs: it assigns high probability to wrong answers it has internalized as facts. These are not epistemic errors (the model is uncertain) but ontological errors (the model's knowledge is wrong). Uncertainty metrics cannot detect what the model does not doubt.

\subsection{Cross-Scale Analysis}

We compare detection performance across two model scales.

\begin{table}[h]
\centering
\caption{Detection Performance Across Model Scales}
\label{tab:scale}
\begin{tabular}{lcccc}
\toprule
\textbf{Model} & \textbf{Hidden Dim} & \textbf{AUC} & \textbf{Cohen's $d$} & \textbf{$p$-value} \\
\midrule
Qwen2.5-1.5B-Instruct & 1536 & 0.532 & $-$0.04 & 0.92 \\
Qwen2.5-7B-Instruct & 3584 & \textbf{0.816} & \textbf{1.06} & 0.018 \\
\bottomrule
\end{tabular}
\end{table}

At 1.5B scale ($d = 1536$), the signal is indistinguishable from random (AUC = 0.532, $p = 0.92$). Notably, Cohen's $d$ is negative ($-$0.04), indicating the signal direction reverses: at 1.5B, hallucinatory outputs have slightly \textit{lower} mean residual magnitude than correct ones, opposite to the 7B pattern. At 7B ($d = 3584$), performance rises to AUC = 0.816. This pattern is consistent with observations that factual knowledge representations strengthen at larger scales~\cite{wei2022emergent}. However, with only two model sizes evaluated, we cannot claim a phase transition; these results demonstrate a scale \textit{difference} requiring validation on intermediate architectures (3B, 13B, 70B). The two models hallucinated on different subsets of the 30 questions (12 vs.\ 11 hallucinations for 1.5B and 7B, respectively), which is expected since model scale affects knowledge accuracy on individual questions.

\subsection{Per-Constraint Ablation}

We decompose the constraint residual at 1.5B scale, where per-token constraint states are available for all five dimensions.

\begin{table}[h]
\centering
\caption{Per-Constraint Detection Performance (Qwen2.5-1.5B-Instruct, $n = 30$)}
\label{tab:components}
\begin{tabular}{lcc}
\toprule
\textbf{Constraint} & \textbf{AUC} & \textbf{Cohen's $d$} \\
\midrule
$\Delta\sigma_{\text{fact}}$ & 0.634 & 0.30 \\
$\Delta\sigma_{\text{style}}$ & 0.588 & 0.10 \\
$\Delta\sigma_{\text{syntax}}$ & 0.472 & $-$0.04 \\
$\Delta\sigma_{\text{safety}}$ & 0.458 & $-$0.10 \\
$\Delta\sigma_{\text{coherence}}$ & 0.431 & $-$0.14 \\
\midrule
$\Delta\|\Pi\|$ (Combined, 1.5B) & 0.532 & $-$0.04 \\
\bottomrule
\end{tabular}
\end{table}

At 1.5B, the factual consistency constraint ($\Delta\sigma_{\text{fact}}$, AUC = 0.634) is the strongest individual predictor, notably \textit{outperforming} the combined 1.5B constraint residual (AUC = 0.532). This suggests that at smaller scales, gradient aggregation across constraints introduces noise that masks the factual signal. The stylistic coherence constraint ($\Delta\sigma_{\text{style}}$, AUC = 0.588) provides a moderate secondary signal.

At 7B, the combined constraint field reaches AUC = 0.816---a substantial improvement over the 1.5B combined result (0.532) and exceeding any individual constraint measured at 1.5B, supporting the hypothesis that larger models develop richer constraint representations where multi-constraint aggregation becomes beneficial. Per-constraint decomposition at 7B scale requires storing all intermediate constraint states during generation and is deferred to future work.

\subsection{Detection Latency}
\label{sec:latency}

Our method uses three model passes: (1) an input encoding pass to extract hidden states from the prompt ($\sim$18.5ms, using \texttt{output\_hidden\_states=True}), (2) the standard autoregressive generation pass ($\sim$1,047ms), and (3) an output re-encoding pass to extract hidden states from the generated text alone ($\sim$10ms), plus constraint computation ($\sim$1ms). Total additional overhead beyond standard generation is $\sim$28.5ms (2.7\% of the 1,076ms total). This is comparable to single-inference baselines (predictive entropy: 1,005ms, which requires only logits from the generation pass) and substantially faster than multi-sample approaches (SelfCheckGPT: 3,011ms for 3$\times$ sampling). The input and output encoding passes operate on short sequences ($<$256 tokens) and add negligible cost relative to autoregressive generation.

\section{Discussion}

\subsection{Comparison to Related Work}

\textbf{Uncertainty-based methods}~\cite{kadavath2022language,kuhn2023semantic} estimate output confidence. They fail when the model is confident in its errors---precisely the regime of false-premise questions.

\textbf{Consistency-based methods}~\cite{manakul2023selfcheckgpt} compare multiple outputs, assuming errors vary while truths are stable. On false-premise questions, this inverts: correct outputs vary (the model reasons through a false premise) while incorrect outputs are consistent (the model reliably recites its misconception).

\textbf{Probe-based methods}~\cite{burns2022discovering,azaria2023internal} train classifiers on hidden states. Our approach shares the intuition that hidden states encode truth-relevant information, but uses \textit{dynamic constraint gradients} across the generation process rather than static probes. This captures the process of constraint violation as it unfolds.

Commercial hallucination detection (Galileo, Arthur, NeMo Guardrails) relies on uncertainty or rules. Constraint residual detection is complementary: it captures a geometric signal invisible to surface-level methods.

\subsection{Limitations}

\textbf{Sample size}: Evaluation used $n = 30$ false-premise questions (11 positive, 19 negative). While the effect size is large (Cohen's $d$ = 1.06, $p = 0.018$), the small sample means AUC variance is itself uncertain. Full-benchmark evaluation on all 817 TruthfulQA questions, with proper confidence intervals via bootstrap, is needed to establish reliable performance bounds.

\textbf{Calibration scale}: Logistic regression on 30 samples (in-sample fit, no regularization) yields preliminary $\alpha=-1.30, \beta=16.52$ estimates. Two-fold cross-validation indicates substantial estimation variance (fold-level AUC range: 0.72--0.80). Expected Calibration Error (ECE) is not reliably estimable at this sample size. Production deployment should use larger, multi-domain calibration sets.

\textbf{Labeling}: Hallucination ground-truth labels are obtained via model self-judgment (the same model evaluates its own outputs). This introduces potential circularity: the model that generates hallucinations also decides what counts as one. Future work should validate against human-annotated labels or multi-model jury judgments.

\textbf{Threshold tuning}: The precision and recall reported in Table~\ref{tab:benchmark} reflect a decision threshold optimized on the evaluation set. Out-of-sample precision and recall at the same threshold will likely be lower. The AUC metric (0.816), being threshold-independent, is the more reliable performance indicator.

\textbf{Model specificity}: Calibration parameters ($\alpha$, $\beta$) and truth/refusal directions are specific to Qwen2.5-7B-Instruct. Cross-model generalization has not been evaluated.

\textbf{Language}: English TruthfulQA only. Cross-lingual performance is unexplored.

\textbf{Latency}: $\sim$1.1s is suitable for post-hoc screening but not for interactive-speed streaming.

\textbf{Scale coverage}: Two model sizes (1.5B, 7B). Intermediate scales and other architectures (Llama, Mistral) are needed to characterize scaling behavior.

\subsection{Future Work}

\begin{itemize}
  \item \textbf{Full-benchmark evaluation}: TruthfulQA 817 questions with bootstrap confidence intervals.
  \item \textbf{Multi-model calibration}: Llama-3, Mistral, DeepSeek architectures; cross-model transfer of constraint directions.
  \item \textbf{Token-level streaming}: Real-time $\Pi(t)$ computation during autoregressive generation for early interception.
  \item \textbf{Cross-lingual validation}: Chinese TruthfulQA, multilingual MMLU.
  \item \textbf{RAG integration}: Compound detection combining constraint residual signals with retrieval grounding.
  \item \textbf{Scaling characterization}: Intermediate scales (3B, 8B, 13B, 70B) to map the constraint field's scaling behavior.
\end{itemize}

\section{Conclusion}

We introduced Constraint Residual Detection, a method that detects hallucinations by measuring violations of a model's internal constraint field rather than output uncertainty. On TruthfulQA false-premise questions with Qwen2.5-7B-Instruct, our method achieves AUC = 0.816 with perfect recall (1.000), substantially outperforming uncertainty-based baselines (best competitor AUC = 0.465, Cohen's $d$ = 1.06, $p$ = 0.018). At 1.5B scale, the signal degrades to near-random (AUC = 0.532), with only the factual consistency constraint retaining discriminatory power (AUC = 0.634), suggesting that multi-constraint aggregation requires sufficient representational capacity.

The key insight is that LLMs \textit{believe} their hallucinations on false-premise tasks---surface confidence is orthogonal to correctness. By observing internal constraint dynamics rather than output distributions, we detect the geometric signature of fabrication that uncertainty-based methods are blind to. The method adds modest computational overhead ($\sim$28.5ms beyond standard generation) and produces monotonic probability estimates via logistic regression, making it practical for hallucination screening when integrated with the generation pipeline.

\begin{thebibliography}{10}

\bibitem{lin2022truthfulqa}
S. Lin, J. Hilton, and O. Evans.
\newblock TruthfulQA: Measuring How Models Mimic Human Falsehoods.
\newblock In \textit{Proceedings of ACL}, 2022.

\bibitem{kadavath2022language}
S. Kadavath, T. Conerly, A. Askell, et al.
\newblock Language Models (Mostly) Know What They Know.
\newblock \textit{arXiv:2207.05221}, 2022.

\bibitem{manakul2023selfcheckgpt}
P. Manakul, A. Liusie, and M. J. F. Gales.
\newblock SelfCheckGPT: Zero-Resource Black-Box Hallucination Detection for Generative Large Language Models.
\newblock In \textit{Proceedings of EMNLP}, 2023.

\bibitem{kuhn2023semantic}
L. Kuhn, Y. Gal, and S. Farquhar.
\newblock Semantic Uncertainty: Linguistic Invariances for Uncertainty Estimation in Natural Language Generation.
\newblock In \textit{Proceedings of ICLR}, 2023.

\bibitem{burns2022discovering}
C. Burns, H. Ye, D. Klein, and J. Steinhardt.
\newblock Discovering Latent Knowledge in Language Models Without Supervision.
\newblock In \textit{Proceedings of ICLR}, 2023.

\bibitem{azaria2023internal}
A. Azaria and T. Mitchell.
\newblock The Internal State of an LLM Knows When It's Lying.
\newblock In \textit{Findings of EMNLP}, 2023.

\bibitem{wei2022emergent}
J. Wei, Y. Tay, R. Bommasani, et al.
\newblock Emergent Abilities of Large Language Models.
\newblock \textit{Transactions on Machine Learning Research}, 2022.

\bibitem{ouyang2022training}
L. Ouyang, J. Wu, X. Jiang, et al.
\newblock Training Language Models to Follow Instructions with Human Feedback.
\newblock In \textit{Proceedings of NeurIPS}, 2022.

\end{thebibliography}

\end{document}
